%! Solving by Square Roots & Completing the Square — Unit 2, Lesson 2.3 — Exit Ticket (Answer Key)
\documentclass[10pt]{article}
\usepackage{algebra2-article}
\usepackage{algebra2-key}   % pulls in -boxes; do NOT also load -boxes
\newcommand{\TallMath}[1]{$\displaystyle #1\rule[-1.4em]{0pt}{3.2em}$}

\begin{document}

\pageheader{Unit 2, Lesson 2.3}{Exit Ticket --- Answer Key}

\vspace{0.12in}

No notes. Show your steps; keep both roots in simplest radical form.

\begin{enumerate}[leftmargin=1.8em, itemsep=12pt]
  \item \textbf{Square Root Property.} Solve $(x-2)^2=18$. Isolate is done for you --- root and simplify.
        \ansline{$x-2=\pm\sqrt{18}=\pm3\sqrt2$, so $x=2\pm3\sqrt2$ ($\approx6.24$ and $-2.24$).}
  \item \textbf{Complete the square.} Solve $x^2+4x-6=0$ by completing the square. Then say what your two
        roots represent on the graph of $y=x^2+4x-6$.
        \ansline{$x^2+4x=6 \Rightarrow x^2+4x+4=10 \Rightarrow (x+2)^2=10 \Rightarrow x=-2\pm\sqrt{10}$.}
        \ansline{The two roots are the $x$-intercepts of the parabola (where $y=0$).}
  \item \textbf{SOL-style multiple choice.} What constant makes $x^2-12x+\underline{\hphantom{00}}$ a
        perfect-square trinomial?
        \begin{enumerate}[label=\Alph*., leftmargin=1.9em, itemsep=1pt]
          \item $6$
          \item $-36$
          \item $36$ \textcolor{keyred}{\textbf{$\leftarrow$ correct}}
          \item $144$
        \end{enumerate}
        \begin{itemize}[leftmargin=1.4em, nosep]
          \item[] Add $\left(\tfrac{b}{2}\right)^2=\left(\tfrac{-12}{2}\right)^2=(-6)^2=36$, giving $x^2-12x+36=(x-6)^2$. \textbf{A} uses $b/2$ (not squared); \textbf{D} uses $b^2$ (forgot to halve); \textbf{B} has the wrong sign.
        \end{itemize}
\end{enumerate}

\end{document}
